\section{Conclusion}
\label{ch:conclusion}

To sum up, what this paper proves is: once the values of context rules are solidified by a corpus and frozen at inference time, no fixed-memory component, whether in the form of a neural network or a Bayesian network, possesses the completeness to independently complete open inversion tasks, and the two theorems share one and the same kernel. The argument unfolds along a decision axis: ill-posedness forces inference to use context, Lemma 1 excludes the global scalar from realizing a heterogeneous profile, Lemma 2 excludes frozen extrinsic coefficients from realizing a heterogeneous profile and converges the criterion onto the source, Criterion P4 and Lemma 3 require values to be instantiated by current evidence, Propositions 4 and 6 complete the classification, and Theorems 5 and 7 close the case. An equivalent, more prudent phrasing holds as well: no frozen rule family can be uniformly optimal on the open set; a frozen family offers at most in-distribution guarantees or guarantees vetted by analytic data consistency, while worst-case guarantees in the pointwise or expected sense require evidence instantiation. The relation to the existing literature has been clarified: the amortization gap records the empirical form of the Mode B gap, No-Free-Lunch gives the algorithm-level universal negation, and this paper fills in the component-level completeness negation between the two and gives a unified decision axis; in relation to No-Free-Lunch, this paper is an instantiation of its converse on the class of open inversion tasks.

This reach is not limited to the two kinds of networks: Markov random fields, conditional random fields, Gaussian mixture models, dictionary learning, kernel methods, and other components, as long as their context rules are carried by parameters trained offline or set by hand, all belong to Mode B and fall within the same falsification; random-field implementations whose coupling strengths are computed online from current evidence belong to Mode A and are outside the negation. The boundary of the conclusion is equally sharp: tasks whose demand can be exactly identified from a fixed corpus are outside the class, and there a fixed-memory component can be complete; the falsification stops exactly where the demand becomes freezable. The exclusion itself is not tied to the quadratic smoothing form of that classical problem: Lemma 1 is proved for a separable data term with an arbitrary even differentiable pairwise potential and without convexity, and its minimal obstruction is an equivalence rather than a sufficient condition, namely one single position whose required strength ratio changes between two admissible configurations. That condition excludes every frozen coefficient family, however many entries the table has; the global scalar is the case in which the ratio must also be common to all positions, which is why the scalar is the shallowest member of the excluded class rather than a special case to be defended separately. For Bayesian networks, beyond the negation there are two further walls of different status: the absence of a perceptual proposal layer, which is the negation read at the input interface rather than an independent theorem, and treewidth blow-up when the source of structure is missing, which is an external hardness result orthogonal to the negation.

Negation is not disparagement: the achievements of these components on closed benchmarks are genuine, and what this paper negates is only the universal claim ``solo use is complete'', not the usefulness of these components in combinations. The algebraic skeleton of the argument chain has been machine-verified with Lean 4 and Mathlib, with zero \texttt{sorry} and zero warnings. The other side of the decision axis, namely the engineering implementation of Mode A, is left to future work.
